\chapter{The ways in}
\label{ch:ways}

\nodwyn{} has one core and several faces. They all speak to the same
sessions, through the same guard, into the same logs, so you can start a
task in one and answer it from another.

\section{The daemon underneath}

\code{nodwyn daemon} is a small server that owns sessions and speaks a
JSON-RPC protocol over a private Unix socket. The browser console, the
terminal client and supervised runs are all clients of it. You rarely start
it by hand: the app starts one inside itself when none is running.

\begin{center}
\begin{tabularx}{\linewidth}{lX}
\toprule
\textbf{command} & \textbf{what you get} \\
\midrule
\code{nodwyn} & a line-mode session in the current folder (no daemon) \\
\code{nodwyn app} & a window (or your browser) with the console; attaches to a running daemon or runs its own \\
\code{nodwyn ui} & the console in your browser, attached to a running daemon \\
\code{nodwyn tui} & a full-screen terminal client, local or attached to a daemon \\
\code{nodwyn daemon} & the session server in the foreground \\
\code{nodwyn run \ldots} & submit a bounded, supervised task to the daemon \\
\code{nodwyn repl}, \code{nodwyn eval} & the Scheme interpreter on its own, no session \\
\bottomrule
\end{tabularx}
\end{center}

\section{The console}

The console, in the app's window or a browser tab, is organised around
tasks. \textbf{Home} starts a task in a folder and lists what needs your
answer, what is running and what finished. A \textbf{task} shows its
conversation with the approval inline (the diff, the reasons, the buttons
from Chapter~\ref{ch:first}), and beside it the plan, this turn's budget,
the changes it made, the project notes and the workspace.
\textbf{Connections} lists the tools (MCP servers among them), the
credential bindings a task loaded, the named workspaces and the models;
\textbf{Workspaces} edits your named folders; and \textbf{Settings} holds
how Home starts a task (a chat, a run or an unattended run), each turn's
limits and the spend limit, the model and its tuning, the judge, the
context window, the keychain and the paired devices. A change applies to
tasks started afterwards; open ones keep what they started with.

\begin{tip}
The console never invents state. Budgets are counted from the log, the plan
is the recorded plan, ``changes'' are recorded writes. If a panel has no
source for something, it says so instead of guessing.
\end{tip}

\section{The terminal client}

\code{nodwyn tui} shows the same conversation full-screen, with a multi-line
composer (\key{Alt+Enter} or \key{Ctrl+J} for a newline, \key{Enter} to
send), and inspectors behind colon commands: \code{:plan}, \code{:budget},
\code{:tools NAME} (a tool's input schema), \code{:bindings},
\code{:replay}, \code{:log [n] [kind]}. \key{Ctrl+C} aborts busy work,
\key{Ctrl+Q} detaches, \key{Ctrl+R} reconnects.

\section{Supervised runs}

A run is a task that goes on for several turns without you typing each
message:

\begin{sh}
nodwyn run --task task.md --repo . --max-turns 20 --budget-wall 2h
nodwyn run status SESSION
nodwyn run stop SESSION
nodwyn run resume SESSION
\end{sh}

Every turn of a run must record a plan before anything irreversible
happens, and between turns you get a \textbf{run checkpoint}: a yes
continues, and grants nothing else. The run ends when the model's final line
is exactly \code{NODWYN\_RUN\_COMPLETE}, which is the model's claim, not a
proof; or when the turn or time limit runs out. With
\code{-{}-unattended} nobody is asked anything, and the run goes on only
after a turn that marked a step of its plan done; with \code{-{}-judge} the
judge answers its ordinary tool confirms. Chapter~\ref{ch:plans} shows the
Scheme that makes this work.

\section{From another device}

The console can also serve other machines, typically your phone over a
tailnet (Tailscale or similar):

\begin{sh}
nodwyn ui --listen 100.101.102.103:7777      # this machine's tailnet address
nodwyn devices                               # who is paired
nodwyn devices revoke phone
\end{sh}

The phone opens the address, enters the pairing code the command printed
\emph{and a device name}, and gets a cookie that keeps it paired for 180
days. Only a hash of its token is stored. Every answer from it is recorded
with its name, \code{via web:phone}, so the log says which screen said yes.

There is no relay in between: the phone talks straight to your computer, so
your computer must be on and reachable, and the network path (a tailnet, or
your home Wi-Fi) is yours to provide.

\begin{careful}
A paired device can answer approvals, and approving a command \emph{is}
running a command as you. Pair only devices you would hand a terminal. A
remote device can never change settings, keys, policy files or the
workspace registry: those need a console on the computer itself.
\end{careful}

\section{Sharing with people who will never open a terminal}

\file{Nodwyn.app} is built for that. It starts its own daemon, pairs its own
window with a one-time ticket in the address's fragment (never sent to a
server, valid for a minute), and walks a first-time user through choosing a
model and creating a keychain. Each person's tasks, keys and settings stay
in their own account. \file{docs/app.md} has the steps for sending it to
someone.
